\section{HERRAMIENTAS Y TÉCNICAS}

El diseño e implementación del sistema demandó la integración de un ecosistema tecnológico multidisciplinario, articulando librerías avanzadas de computación científica, marcos de aprendizaje automático, tecnologías geoespaciales y arquitecturas cliente-servidor para visualización en tiempo real:

\begin{itemize}
    \item \textbf{Lenguaje de Programación y Entorno de Ejecución:} Python 3.12 como lenguaje central de procesamiento analítico y backend web, operando bajo un entorno virtual aislado (`venv`) en Windows 11.
    \item \textbf{Computación Numérica y Manipulación de Datos:}
    \begin{itemize}
        \item \textbf{NumPy:} Álgebra lineal de vectores densos, transformaciones trigonométricas angulares $(\sin, \cos)$, conversión a radianes y funciones de cálculo métrico euclidiano y geodésico.
        \item \textbf{Pandas y PyArrow:} Indexación, limpieza, filtrado relacional y serialización columnar en formato Parquet comprimido con algoritmo Snappy para optimización de I/O en disco.
    \end{itemize}
    \item \textbf{Ecosistema de Machine Learning y Modelado No Supervisado:}
    \begin{itemize}
        \item \textbf{scikit-learn:} Implementación de KMeans (con inicialización K-Means++), DBSCAN (con métrica `haversine`), OPTICS, AgglomerativeClustering (enlace de Ward), GaussianMixture, PCA, TruncatedSVD, FactorAnalysis, TSNE, TfidfVectorizer, CountVectorizer, LatentDirichletAllocation, y MLPRegressor para el Autoencoder de anomalías.
        \item \textbf{hdbscan:} Paquete especializado optimizado en Cython para clustering jerárquico por exceso de masa y estabilidad persistente $\lambda$.
        \item \textbf{umap-learn:} Librería de reducción no lineal basada en geometría riemanniana y grafos difusos para proyecciones de alta fidelidad.
        \item \textbf{SciPy:} Descomposición espectral, cálculo de distancias por enlace, matrices de condensación y coeficientes de correlación cofenética.
    \end{itemize}
    \item \textbf{Framework de Servicio Web y API REST:}
    \begin{itemize}
        \item \textbf{Flask:} Microframework web en Python configurado mediante arquitectura modular con Blueprints (`modules/routes/`), manejador de errores y endpoints JSON RESTful de baja latencia.
        \item \textbf{Jinja2:} Motor de renderizado de plantillas HTML5 reactivas del lado del servidor.
        \item \textbf{Joblib:} Serialización y deserialización binaria de modelos pre-entrenados para inferencia en tiempo real en memoria.
    \end{itemize}
    \item \textbf{Visualización Cartográfica y Gráficos Interactivos:}
    \begin{itemize}
        \item \textbf{MapLibre GL JS / Deck.gl / Three.js:} Motor gráfico cliente acelerado por GPU (WebGL) para el renderizado del visor esférico inmersivo (Globo 3D), proyecciones esféricas, capas de dispersión de crímenes y extrusiones 3D.
        \item \textbf{Folium:} Generación de mapas interactivos basados en Leaflet para visualización 2D de capas HeatMap y agrupamientos espaciales.
        \item \textbf{Matplotlib y Seaborn:} Generación estática de alta resolución de dendrogramas, diagramas de caja (boxplots), matrices de correlación y curvas de silueta.
    \end{itemize}
    \item \textbf{Control de Versiones y Documentación Técnica:} Git y GitHub para control de versiones distribuido, seguimiento de código modular y trazabilidad de experimentos.
\end{itemize}

